\section{Limits, nonclaims, and research frontier}

A useful framework should make its limits more precise than its aspirations.

\subsection{The normal form does not prove every theorem in iPRA}

\(\FQNF\) is a compilation theorem.  It says that once a theorem's finite
producer, verifier, and extractor are available in object iPRA, finite evidence
can be normalized and selected.  It does not imply that every classical
mathematical theorem has such a producer, or that iPRA proves the producer's
existence.

In particular, the framework does not establish general effective forms of:

\begin{itemize}
\item Faltings' or Siegel's theorem;
\item Mordell--Weil representative production;
\item prime-number-theorem constants or zero-free regions;
\item unrestricted compactness or choice;
\item exact selectors at discontinuities;
\item a universal PR evaluator;
\item arbitrary higher-type comprehension.
\end{itemize}

It gives a rigorous conditional endpoint for each: an explicit producer field
whose satisfaction is sufficient for the finite remainder of the proof.

\subsection{Brute-force search may be enormous}

Whole-code search is primitive recursive but can be exponentially or worse
than specialized algorithms.  It should be treated as a normalization and
existence compiler.  After formalization, efficient producers can be proved
accepted by the same checker without changing downstream theorems.

There is also a proof-size tradeoff.  A monolithic global checker can be easier
to state but harder to reason about.  The cut-based modular interface remains
important for human proofs and Lean elaboration.

\subsection{A numerical bound is indispensable}

The framework cannot search a merely finite set unless its codes lie below an
explicit term.  In some applications the code bound itself may be huge but
elementary; in others its production is equivalent to the main deep theorem.
The boundary is mathematical, not an artifact of this encoding.

\subsection{Reflection is real work}

Writing a Python verifier or Lean Boolean function is not enough.  The object
checker must be represented by fixed PR syntax, and its soundness must be a
\(\Prf\,\PRAAx\) derivation.  The report reduces this work to reusable Boolean,
digit, fold, and search lemmas, but those lemmas still have to be checked.

\subsection{Coding choices affect engineering}

Uniform base, mixed radix, nested pairing, prime-power coding, and Chinese
remainder coding are all possible.  Uniform base has the simplest fixed-field
semantics; mixed radix gives tighter bounds; pairing reuses a mature checked
module.  The first formal integration should choose whichever minimizes proof
engineering, not whichever gives the shortest numerical code.

\subsection{Metatheoretic classical axioms}

The existing Lean developments may use ordinary classical principles in the
metatheory while proving an intuitionistic object theorem.  The meaningful
audit question is whether the final declaration's object assumptions exceed
\(\PRAAx\), not whether Lean's implementation language is intuitionistic.
Nevertheless, every public theorem should have its Lean axioms printed and
compared with the repository's accepted audit baseline.

\subsection{The immediate research frontier}

The highest-value next results are concrete and modest in number.

\begin{enumerate}
\item Kernel-check the structural core.
\item Formalize mixed-radix coding and digit recovery in
\code{IPRADerivations}.
\item Prove generic reflected bounded folds.
\item Complete least-certificate selection from existential production.
\item Formalize the three-level path benchmark.
\item Formalize rational bisection and interval stability.
\item Refactor existing packets to instantiate the package rather than creating
new parallel frameworks.
\end{enumerate}

After these steps, many proof packets can be classified cleanly: complete
instance, missing local arithmetic, missing quantitative producer, or
constructively obstructed exact statement.
